\section{Related Work}
\label{sec:related}

\subsection{Benchmark Contamination Detection}
\label{sec:related:detection}

Detecting training data contamination in LLMs has produced a diverse methodological landscape spanning at least eight paradigms. Time-based partitioning uses temporal cutoffs to identify potentially leaked items, while canary string insertion (e.g., BIG-Bench) provides direct memorization probes \citep{chang2025benchmarking}. Min-K\%++ \citep{duan2025minkpp} performs membership inference by examining token-level probability distributions, achieving strong detection on pre-training contamination. Time Travel \citep{oren2024proving} uses guided completion to detect black-box contamination without model weight access. ConStat \citep{dekoninck2024constat} measures performance gaps between original and paraphrased benchmark items, targeting surface-form memorization. Self-Critique \citep{wust2025selfcritique} identifies contamination through entropy collapse in model self-evaluations---a signal that, unlike probability-based methods, partially survives reinforcement learning post-training. Hierarchical semantic detection approaches match items at multiple granularity levels to training corpora. KDS \citep{xu2025kds} uses knowledge distillation signatures to achieve near-perfect detection (Spearman $\rho = 0.99$) but requires per-model LoRA fine-tuning, limiting scalability.

A critical finding is that detection methods are fragile under RL post-training: \citet{wang2025fragility} show that GRPO and similar procedures degrade probability-based detection signals substantially, while Self-Critique's entropy-based signal achieves only 0.60--0.70 AUC on RL-trained models \citep{wust2025selfcritique}. Despite this rich detection toolkit, all existing methods produce binary or continuous \emph{flags}---they diagnose contamination without prescribing how to adjust the reported accuracy that end-users rely on.

\subsection{Contamination Impact and Mitigation}
\label{sec:related:mitigation}

Understanding contamination's quantitative impact has advanced rapidly. \citet{dong2026dose} establish dose-response curves showing that score inflation scales monotonically with both model size and contamination dosage. \citet{yang2026rl} further demonstrate that SFT inflates accuracy only on contaminated tasks, while GRPO amplifies inflation even on structurally similar uncontaminated items---a cross-contamination effect that complicates attribution. \citet{zhang2025forget} show that contamination persists in model weights following standard decontamination, with forgetting curves following scaling laws that make complete removal impractical.

The mitigation landscape is equally sobering. \citet{xu2025emperor} systematically evaluate twenty strategies---item rephrasing, option shuffling, filtering, dynamic benchmark construction, and more---across ten LLMs and five benchmarks, finding that none reliably eliminates contamination effects. Dynamic benchmarks such as LiveBench \citep{white2024livebench} circumvent contamination by generating fresh items but cannot retroactively calibrate existing benchmark scores. Inference-Time Decontamination \citep{yang2024inference} rewrites items at test time but operates during inference rather than producing post-hoc calibrated scores. DCR \citep{li2025dcr} is the closest existing method to ours, producing adjusted accuracy via fuzzy inference over detection signals; however, it provides point estimates without uncertainty bounds and does not address signal degradation under RL. \citet{rao2026silicon} propose a router-worker audit framework that produces contamination flags rather than calibrated scores. Collectively, these findings motivate a paradigm shift: from contamination \emph{prevention}---which demonstrably fails---to post-hoc statistical \emph{adjustment} with proper uncertainty quantification.

\subsection{IRT for LLM Evaluation}
\label{sec:related:irt}

Item response theory has gained traction as a principled framework for LLM evaluation. \citet{lalor2024irt} demonstrate that IRT-based evaluation yields more informative model comparisons than raw accuracy. PSN-IRT \citep{gu2026psnirt} scales this to 41,871 items across 11 benchmarks using a 4PL model with neural network parameter estimation, establishing item difficulty and discrimination profiles but without modeling contamination or producing adjusted scores. \citet{polo2026growing} fit multidimensional 2PL IRT with anchor items to over 400 models, validating IRT's applicability at scale while identifying practical limitations in the few-examinee regime; contamination is not addressed. ATLAS \citep{lyu2025atlas} leverages IRT for adaptive testing, reducing evaluation cost through maximally informative item selection. TinyBenchmarks \citep{polo2024tinybenchmarks} similarly exploits IRT-informed item selection for efficient evaluation. HiBayES \citep{zhuang2025hibayes} provides hierarchical Bayesian credible intervals over model ability estimates but treats all items as equally trustworthy. \citet{han2026probing} quantify contamination effects within psychometric evaluations but do not use IRT or produce calibrated scores. Across this body of work, contamination has not been modeled as a covariate within the IRT framework---a gap that ContamCal addresses by augmenting the item response function with exposure parameters.

\subsection{Item Compromise in Educational Measurement}
\label{sec:related:edumeasure}

The psychometric literature has long addressed item compromise---the unauthorized disclosure of test content prior to administration. \citet{eckerly2017} develops IRT models that incorporate compromise indicators to simultaneously estimate item parameters and compromise effects, while \citet{sinharay2017} proposes likelihood ratio and score tests for detecting item preknowledge. The explanatory IRT tradition \citep{deboeck2004,adams1997} provides the theoretical foundation for embedding external covariates within item response functions, and \citet{mislevy1986} establishes Bayesian estimation procedures for such augmented models.

ContamCal transfers this item compromise methodology to LLM evaluation through three non-trivial structural adaptations: (a)~independent examinees become a single model producing all responses through shared parameters, fundamentally altering local independence semantics; (b)~binary leakage indicators are replaced with continuous exposure scores derived from noisy detection signals, necessitating a measurement error model \citep{maris2009} that propagates detection uncertainty into ability estimates; and (c)~stable examinee ability gives way to regime-conditional behavior, as SFT and RL post-training produce qualitatively different contamination signatures affecting both detection reliability and effect magnitude.

\paragraph{Positioning.} ContamCal bridges three previously disconnected research threads: contamination detection (which flags but does not calibrate), IRT-based LLM evaluation (which models ability but ignores contamination), and educational item compromise models (which handle leaked items but assume independent examinees with binary exposure). By augmenting a 4PL IRT model with multi-signal exposure covariates and a measurement error model, ContamCal produces calibrated benchmark scores with Bayesian credible intervals---providing uncertainty quantification that both detection-only methods and prior calibration attempts lack, while remaining robust to the signal degradation induced by RL post-training.
